% !TEX root = ../main.tex
\corrigesChapitre{20}{Lumière : ondes et photons}
\corrigesSouples

\begin{blocLivre}{cor-20-1}
\begin{corrige}{1}{Des ondes au quotidien}
\demarche{Convertir les GHz en Hz et les nm ou \unit{\micro m} en m, puis utiliser $\lambda=\dfrac{c}{\nu}$ ou $\nu=\dfrac{c}{\lambda}$.}
\begin{questions}
\item (a) $\lambda=\dfrac{\num{3,00e8}}{\num{2,4e9}}=\res{\qty{0,125}{m}}$, soit \qty{12,5}{cm} ; (b) $\lambda=\dfrac{\num{3,00e8}}{\num{2,45e9}}=\res{\qty{0,122}{m}}$, soit \qty{12,2}{cm}.
\item (c) $\nu=\dfrac{\num{3,00e8}}{\num{0,10e-9}}=\res{\qty{3,0e18}{Hz}}$ ; (d) $\nu=\dfrac{\num{3,00e8}}{\num{10e-6}}=\res{\qty{3,0e13}{Hz}}$.
\item (a) et (b) : \res{\text{micro-ondes}} ; (c) : \res{\text{rayons X}} ; (d) : \res{\text{infrarouge}}.
\item Plus grande fréquence : \res{\text{(c)}}, les rayons X ; plus grande longueur d'onde : \res{\text{(a)}}, le wifi.
\end{questions}
\schema[corrige]{20-ex01-cor}
\piege{Oublier de convertir : \qty{1}{GHz}~$=\qty{e9}{Hz}$, \qty{1}{nm}~$=\qty{e-9}{m}$, \qty{1}{\micro m}~$=\qty{e-6}{m}$.}
\end{corrige}
\end{blocLivre}
\begin{blocLivre}{cor-20-2}
\begin{corrige}{2}{Photons rouge et violet}
\demarche{$E=\dfrac{h\,c}{\lambda}$ avec $\lambda$ en mètres donne $E$ en joules ; on divise par \num{1,60e-19} pour l'avoir en eV.}
\begin{questions}
\item $E=\dfrac{\num{6,63e-34}\times\num{3,00e8}}{\num{700e-9}}=\res{\qty{2,84e-19}{J}}$, soit $\dfrac{\num{2,84e-19}}{\num{1,60e-19}}=\res{\qty{1,78}{eV}}$.
\item $E=\dfrac{\num{6,63e-34}\times\num{3,00e8}}{\num{400e-9}}=\res{\qty{4,97e-19}{J}}$, soit \res{\qty{3,11}{eV}}.
\item Le photon \res{\text{violet}}, $\dfrac{700}{400}=\res{1{,}75}$ fois plus énergétique : l'énergie est inversement proportionnelle à $\lambda$.
\item Par énergie croissante, donc longueur d'onde décroissante : \res{\text{infrarouge, visible, ultraviolet, X}}.
\end{questions}
\retiens{Plus la longueur d'onde est courte, plus le photon est énergétique : le violet transporte presque deux fois plus d'énergie que le rouge.}
\end{corrige}
\end{blocLivre}
\begin{blocLivre}{cor-20-3}
\begin{corrige}{3}{Le diagramme de l'hydrogène}
\demarche{Une transition vers le bas émet un photon d'énergie $\Delta E=E_\text{haut}-E_\text{bas}$ ; on en déduit $\lambda=\dfrac{h\,c}{\Delta E}$, avec $\Delta E$ en joules.}
\schemaCote{20-ex03}[0.42]{\begin{questions}
\item État fondamental : le niveau le plus bas, \res{n=1} (\qty{-13,6}{eV}). Le niveau 0 correspond à l'\res{\text{atome ionisé}} : l'électron est arraché.
\item L'atome descend : il \res{\text{émet}} un photon d'énergie $\Delta E=-3{,}40-(-13{,}6)=\res{\qty{10,2}{eV}}$, soit $10{,}2\times\num{1,60e-19}=\res{\qty{1,63e-18}{J}}$.
\item $\lambda=\dfrac{\num{6,63e-34}\times\num{3,00e8}}{\num{1,63e-18}}=\qty{1,22e-7}{m}=\res{\qty{122}{nm}}$ : domaine des \res{\text{ultraviolets}}.
\item $\Delta E=-1{,}51-(-13{,}6)=\qty{12,1}{eV}=\qty{1,93e-18}{J}$, d'où $\lambda=\res{\qty{103}{nm}}$, encore dans l'\res{\text{ultraviolet}}.
\end{questions}}
\end{corrige}
\end{blocLivre}
\begin{blocLivre}{cor-20-4}
\begin{corrige}{4}{La lampe au sodium}
\demarche{Calculer $\nu=\dfrac{c}{\lambda}$, puis $E=h\,\nu$, et comparer à l'écart entre les deux niveaux.}
\begin{questions}
\item $\nu=\dfrac{\num{3,00e8}}{\num{589e-9}}=\res{\qty{5,09e14}{Hz}}$.
\item $E=h\,\nu=\num{6,63e-34}\times\num{5,09e14}=\res{\qty{3,38e-19}{J}}$, soit \res{\qty{2,11}{eV}}.
\item $E_1-E_0=-3{,}03-(-5{,}14)=\qty{2,11}{eV}$ : c'est \res{\text{exactement l'énergie du photon}}, émis quand l'atome passe de $E_1$ à $E_0$.
\item La voiture ne diffuse que le bleu ; la lampe n'émet que du jaune-orangé, que la carrosserie absorbe : elle paraît \res{\text{noire}}.
\end{questions}
\end{corrige}
\end{blocLivre}
\begin{blocLivre}{cor-20-5}
\begin{corrige}{5}{Micro-ondes et ultraviolets}
\demarche{Comparer l'énergie d'un seul photon à l'énergie d'une liaison : une molécule absorbe les photons un par un.}
\begin{questions}
\item $E=h\,\nu=\num{6,63e-34}\times\num{2,45e9}=\res{\qty{1,62e-24}{J}}$, soit \res{\qty{1,0e-5}{eV}}.
\item $E=\dfrac{\num{6,63e-34}\times\num{3,00e8}}{\num{300e-9}}=\qty{6,63e-19}{J}$, soit \res{\qty{4,14}{eV}}.
\item $\dfrac{4{,}14}{\num{1,0e-5}}\approx\res{\num{4e5}}$ : le photon UV est environ \num{400000} fois plus énergétique.
\item Un photon UV apporte à lui seul quelques eV : il \res{\text{peut casser une liaison}} de l'ADN. Un photon micro-onde apporte cent mille fois
trop peu : même très nombreux, ces photons ne font que \res{\text{chauffer}} les molécules.
\end{questions}
\retiens{L'effet d'un rayonnement sur une molécule dépend de l'énergie de chaque photon, pas du nombre de photons reçus.}
\end{corrige}
\end{blocLivre}
\begin{blocLivre}{cor-20-6}
\begin{corrige}{6}{Les raies de l'hydrogène, en Python}
\demarche{Le programme applique $\Delta E=E_\text{haut}-E_\text{bas}$ puis $\lambda=\dfrac{h\,c}{\Delta E}$ ; la fonction \code{domaine} compare $\lambda$ aux limites du visible.}
\begin{questions}
\item $E_\text{haut}-E_\text{bas}=-\dfrac{13{,}6}{n_\text{haut}^2}+\dfrac{13{,}6}{n_\text{bas}^2}$ : c'est \res{\code{dE}}, positif car $n_\text{bas}<n_\text{haut}$, en eV.
\item $\Delta E=13{,}6-\dfrac{13{,}6}{4}=\qty{10,2}{eV}$, $\lambda=\dfrac{\num{6,63e-34}\times\num{3,00e8}}{10{,}2\times\num{1,60e-19}}=\res{\qty{122}{nm}}$.
\item \res{\text{Non}} : selon le niveau d'arrivée, il émet dans l'ultraviolet ($n=1$), le visible ($n=2$) ou l'infrarouge ($n=3$).
\item Ajouter le couple \code{(5, 2)} à la liste ; affichage : \res{\code{5 -> 2 : 435 nm visible}}.
\end{questions}
\end{corrige}
\end{blocLivre}
\begin{blocLivre}{cor-20-7}
\begin{corrige}{7}{Le pointeur laser}
\demarche{Chaque seconde, le laser émet \qty{1,0e-3}{J} : le nombre de photons est le quotient de cette énergie par celle d'un photon.}
\begin{questions}
\item $\nu=\dfrac{\num{3,00e8}}{\num{650e-9}}=\res{\qty{4,62e14}{Hz}}$ ; $E=h\,\nu=\res{\qty{3,06e-19}{J}}$, soit \res{\qty{1,91}{eV}}.
\item $N=\dfrac{\num{1,0e-3}}{\num{3,06e-19}}=\res{\num{3,27e15}}$ photons par seconde.
\item $E=\dfrac{h\,c}{\num{532e-9}}=\qty{3,74e-19}{J}$, d'où $N=\num{2,67e15}$ : \res{\text{moins de photons}}, chacun étant plus énergétique.
\item $E=P\,\Delta t=\num{1,0e-3}\times60=\res{\qty{0,060}{J}}$. C'est peu, mais le cristallin concentre le faisceau sur un minuscule point de la
rétine, qui peut être \res{\text{brûlé}}.
\end{questions}
\end{corrige}
\end{blocLivre}
\begin{blocLivre}{cor-20-8}
\begin{corrige}{8}{La lampe à vapeur de mercure}
\demarche{Pour chaque transition : $\Delta E$ en eV, conversion en joules, puis $\lambda=\dfrac{h\,c}{\Delta E}$.}
\schemaCote{20-ex08}[0.42]{\begin{questions}
\item $E_4\to E_1$ : $\Delta E=\qty{2,84}{eV}$, $\lambda=\dfrac{\num{6,63e-34}\times\num{3,00e8}}{2{,}84\times\num{1,60e-19}}=\res{\qty{438}{nm}}$, \res{\text{bleu}} ; $E_4\to E_2$ : $\Delta E=\qty{2,28}{eV}$, \res{\qty{545}{nm}}, \res{\text{vert}}.
\item $\Delta E=-5{,}55-(-10{,}44)=\qty{4,89}{eV}$, d'où $\lambda=\res{\qty{254}{nm}}$ : \res{\text{ultraviolet}}.
\item Cette raie est invisible et dangereuse : sans la poudre, le tube éclairerait très peu. La poudre \res{\text{absorbe les UV et réémet du visible}}.
\item Il faudrait un niveau à $-10{,}44+4{,}00=\qty{-6,44}{eV}$ : il n'existe pas, le photon de \qty{4,00}{eV} \res{\text{n'est pas absorbé}}. Celui de
\qty{4,89}{eV} l'est (passage à $E_1$).
\end{questions}}
\end{corrige}
\end{blocLivre}
\begin{blocLivre}{cor-20-9}
\begin{corrige}{9}{La radiographie}
\demarche{Mêmes relations que pour la lumière visible ; seule l'échelle change.}
\begin{questions}
\item $\nu=\dfrac{\num{3,00e8}}{\num{0,050e-9}}=\res{\qty{6,0e18}{Hz}}$.
\item $E=\dfrac{h\,c}{\lambda}=\dfrac{\num{6,63e-34}\times\num{3,00e8}}{\num{0,050e-9}}=\res{\qty{3,98e-15}{J}}$, soit \res{\qty{2,49e4}{eV}}.
\item $\dfrac{500}{0{,}050}=\res{\num{1,0e4}}$ : dix mille fois plus énergétique qu'un photon visible.
\item Le détecteur placé derrière le patient n'est atteint que là où les X ont traversé : les os y laissent une \res{\text{ombre}}. Ces photons
très énergétiques \res{\text{ionisent}} les molécules : à chaque cliché, le manipulateur s'abrite derrière une paroi plombée.
\end{questions}
\end{corrige}
\end{blocLivre}
\begin{blocLivre}{cor-20-10}
\begin{corrige}{10}{Absorption par l'hydrogène}
\demarche{Depuis le niveau $n=1$, un photon n'est absorbé que si $E_1+E_\text{photon}$ tombe exactement sur un niveau, ou dépasse 0 (ionisation).}
\schemaCote{20-ex10-cor}[0.42]{\begin{questions}
\item $E_2-E_1=\qty{10,2}{eV}$ et $E_3-E_1=\qty{12,1}{eV}$ : ces deux photons sont \res{\text{absorbés}}. Avec \qty{11,0}{eV}, l'atome arriverait à
\qty{-2,6}{eV}, où il n'y a aucun niveau : \res{\text{pas d'absorption}}.
\item $0-(-13{,}6)=\res{\qty{13,6}{eV}}$ ; $\lambda=\dfrac{\num{6,63e-34}\times\num{3,00e8}}{13{,}6\times\num{1,60e-19}}=\res{\qty{91,4}{nm}}$ (UV).
\item L'atome est \res{\text{ionisé}} ; l'excédent devient l'énergie cinétique de l'électron : $15{,}0-13{,}6=\res{\qty{1,4}{eV}}$.
\item Les atomes sont tous dans l'état fondamental ; depuis $n=1$, le plus petit photon absorbable a \qty{10,2}{eV} (UV). Les photons visibles,
de \qty{1,6}{eV} à \qty{3,1}{eV}, \res{\text{ne correspondent à aucune transition}} depuis ce niveau.
\end{questions}}
\piege{Croire qu'un atome absorbe tout photon assez énergétique : sous le seuil d'ionisation, il faut un écart exact entre deux niveaux.}
\end{corrige}
\end{blocLivre}
\begin{blocLivre}{cor-20-11}
\begin{corrige}{11}{Les DEL de couleur}
\demarche{$\lambda=\dfrac{h\,c}{E}$ avec $E$ converti en joules ; on situe ensuite $\lambda$ dans le spectre visible.}
\begin{questions}
\item $\lambda=\dfrac{\num{6,63e-34}\times\num{3,00e8}}{1{,}9\times\num{1,60e-19}}=\res{\qty{654}{nm}}$ : \res{\text{rouge}}.
\item $\lambda=\dfrac{\num{6,63e-34}\times\num{3,00e8}}{2{,}76\times\num{1,60e-19}}=\res{\qty{450}{nm}}$ : \res{\text{bleu}}.
\item $E=\dfrac{\num{6,63e-34}\times\num{3,00e8}}{\num{525e-9}}=\qty{3,79e-19}{J}$, soit \res{\qty{2,37}{eV}}.
\item Le bleu et le jaune sont des couleurs \res{\text{complémentaires}} : leur superposition donne du blanc (synthèse additive, chapitre 19).
\end{questions}
\end{corrige}
\end{blocLivre}
\begin{blocLivre}{cor-20-12}
\begin{corrige}{12}{Le four à micro-ondes}
\demarche{La distance entre deux zones fondues donne $\dfrac{\lambda}{2}$ ; on la mesure sur plusieurs intervalles pour gagner en précision.}
\begin{questions}
\item $\lambda=\dfrac{\num{3,00e8}}{\num{2,45e9}}=\res{\qty{0,122}{m}}=\qty{12,2}{cm}$, compris entre \qty{1}{mm} et \qty{1}{m} : ce sont bien des \res{\text{micro-ondes}}.
\item Sur deux intervalles : $13{,}2-1{,}0=\qty{12,2}{cm}$, soit $\dfrac{\lambda}{2}=\res{\qty{6,1}{cm}}$ et $\lambda=\qty{12,2}{cm}$. Alors
$c=\lambda\,\nu=0{,}122\times\num{2,45e9}=\res{\qty{2,99e8}{m.s^{-1}}}$, en accord avec la valeur attendue.
\item Un photon : $E=h\,\nu=\qty{1,62e-24}{J}$ ; chaque seconde : $N=\dfrac{800}{\num{1,62e-24}}=\res{\num{4,9e26}}$ photons.
\item Les zones chaudes sont fixes dans le four : en tournant, l'aliment les traverse toutes et \res{\text{chauffe uniformément}}.
\end{questions}
\schema[corrige]{20-ex12}
\end{corrige}
\end{blocLivre}
\begin{blocLivre}{cor-20-13}
\begin{corrige}{13}{Les feux d'artifice}
\demarche{Chaque élément émet les photons de ses propres transitions ; les lumières de plusieurs sels s'ajoutent dans l'œil.}
\begin{questions}
\item Les atomes \res{\text{émettent}} de la lumière en se désexcitant ; les niveaux étant quantifiés, ils n'émettent que certaines
radiations : c'est un \res{\text{spectre de raies}}.
\item Sodium : jaune-orangé, \res{\qty{2,11}{eV}} ; strontium : rouge, \res{\qty{1,91}{eV}} ; baryum : vert, \res{\qty{2,39}{eV}}.
\item Rouge + vert = jaune en synthèse additive : mélanger des sels de \res{\text{strontium et de baryum}}.
\item Le pull ne diffuse que le bleu ; il ne reçoit que du jaune-orangé, qu'il absorbe : il paraît \res{\text{noir}}.
\end{questions}
\end{corrige}
\end{blocLivre}
\begin{blocLivre}{cor-20-14}
\begin{corrige}{14}{La cellule photovoltaïque}
\demarche{Comparer l'énergie de chaque photon au seuil de \qty{1,12}{eV} : en dessous, rien n'est produit ; au-dessus, l'excédent chauffe la cellule.}
\begin{questions}
\item $\lambda_\text{max}=\dfrac{h\,c}{E}=\dfrac{\num{6,63e-34}\times\num{3,00e8}}{1{,}12\times\num{1,60e-19}}=\res{\qty{1,11e-6}{m}}$, soit \qty{1110}{nm} : \res{\text{infrarouge}}.
\item $E=\qty{2,49}{eV}$ ; perdue : $2{,}49-1{,}12=\res{\qty{1,37}{eV}}$, soit $\dfrac{1{,}37}{2{,}49}=\res{\qty{55}{\percent}}$ de son énergie.
\item $E=\dfrac{h\,c}{\num{1500e-9}}=\qty{0,83}{eV}<\qty{1,12}{eV}$ : \res{\text{non}}, il traverse ou chauffe la cellule.
\item Une grande partie de l'infrarouge solaire est \res{\text{inutilisable}}, et chaque photon absorbé \res{\text{perd son excédent}} en chaleur : le
rendement plafonne vers \qty{20}{\percent} pour les panneaux courants.
\end{questions}
\schema[corrige]{20-ex14-cor}
\end{corrige}
\end{blocLivre}
\begin{blocLivre}{cor-20-15}
\begin{corrige}{15}{Les raies de Fraunhofer}
\demarche{Calculer l'énergie des photons manquants et la comparer aux écarts entre niveaux de l'hydrogène.}
\begin{questions}
\item Spectre d'\res{\text{absorption}} : la surface chaude du Soleil émet un spectre continu ; les gaz de son atmosphère \res{\text{absorbent}}
certaines radiations, qui manquent.
\item $E=\dfrac{\num{6,63e-34}\times\num{3,00e8}}{\num{656,3e-9}}=\qty{3,03e-19}{J}=\res{\qty{1,89}{eV}}$ ; à \qty{486,1}{nm} : \qty{4,09e-19}{J}~$=\res{\qty{2,56}{eV}}$.
\item $E_2=\qty{-3,40}{eV}$, $E_3=\qty{-1,51}{eV}$, $E_4=\qty{-0,85}{eV}$ : $E_3-E_2=\qty{1,89}{eV}$ et $E_4-E_2=\qty{2,55}{eV}$, aux arrondis près. Ces
photons sont absorbés par des atomes d'\res{\text{hydrogène}} passant du niveau 2 aux niveaux 3 et 4.
\item La raie à \qty{589,0}{nm} est celle du \res{\text{sodium}} (exercice 4).
\end{questions}
\schema[corrige]{20-ex15}
\end{corrige}
\end{blocLivre}
